\section{A sharp finite-memory law for an observable decision}\label{sec:memory-risk}
Fix $n,m\ge1$, the apparatus, the prior, and a probe list from
Lemma~\ref{lem:probe-span}. The following single decision problem is fixed
before any history is compressed. After $n$ reports the controller stores
its state. An independent public index $J$, uniform on $\{1,\ldots,s\}$,
is then revealed. The forecaster announces $a\in[0,1]$. The $m$ cartridges
of probe $J$ are executed and $Y$ records whether all $m$ reports are
$\dagger$. The terminal reward is $1-(a-Y)^2$. There is no acceptance
cost and no direct reward on the unobserved parameter. Exactly $n+m$
cartridges are consumed, including failed placement and censored records.

Immutable apparatus data, $n,m$ and the newly revealed index $J$ are known
to the decoder. Past commands, past reports and arbitrary history-dependent
real parameters may not be kept outside the charged state. The index is
revealed after encoding; otherwise a single predicted scalar would answer
a different, easier question. This is a one-checkpoint memory problem,
not an unproved assertion about repeated lossy quantization at every time.
The common code is selected before the past command program is supplied.
That program is an input stream, not part of immutable calibration. The
worst-history assertion is uniform over its admitted realizations. Below
we also prove an unconditional version for one fixed randomized prefix
experiment, so a different code for each past program is never used.

\begin{lemma}[The score identifies the operational quotient]\label{lem:score-metric}
Let $p=(p_j)_{j=1}^s$ be a history's probe prediction vector. With full
history the unique Bayes forecast for query $j$ is $p_j$. For any decoded
vector $a=(a_j)$ the excess conditional squared loss is
\begin{equation}\label{eq:score-risk}
 \mathcal R(a,p)=s^{-1}\|a-p\|_2^2.
\end{equation}
The same formula holds with expectation over randomized decoded forecasts.
There exist $p_*$, an isometric linear map $E:\mathbb R^d\to\mathbb R^s$,
and $\rho>0$, $d=q\min(n,m)$, such that all predictions
\begin{equation}\label{eq:prediction-ball}
 p_*+\rho E u,\qquad \|u\|_2\le1,
\end{equation}
are attainable on a compact family of all-failure histories, with a
continuous section. Every point of this ball lies in $[0,1]^s$.
\end{lemma}
\begin{proof}
For $Y\sim\operatorname{Bernoulli}(p_j)$,
$\mathbb E(a_j-Y)^2=p_j(1-p_j)+(a_j-p_j)^2$. Average over $J$.
The closed ball and section follow by shrinking the relatively open
prediction patch in Theorem~\ref{thm:decision-quotient}. Choose an
orthonormal basis of its tangent space to obtain $E$. The ball consists
of actual event probabilities, so its inclusion in the cube is automatic.
\end{proof}

\paragraph{A quantitative choice of the prediction ball.}
The radius can be bounded using finite calibration data, rather than
regarded as an unspecified geometric constant. Choose an interior coprime
factor tuple and a $d$-dimensional linear slice of its gate coordinates on
which the prediction derivative has rank $d$. Let $u=0$ denote that tuple,
choose orthonormal coordinates in the prediction affine space, and write
$\phi(u)=E^\top(p(u)-p(0))$. Set $A=D\phi(0)$ and
$\sigma=\sigma_{\min}(A)>0$. Choose $r>0$ so the closed gate slice
$\|u\|\le r$ stays strictly within the allowed cube and
\[
 \sup_{\|u\|\le r}\|D\phi(u)-A\|\le\sigma/2.
\]
Then one may take $\rho=\sigma r/2$. Indeed for every $\|y\|\le\rho$ the
map $u\mapsto A^{-1}\{y-\phi(u)+Au\}$ sends the radius-$r$ ball into
itself and is a contraction with constant at most $1/2$. Its unique fixed
point satisfies $\phi(u)=y$ and depends continuously on $y$. Since all
predictions lie in the same affine plane, this proves
\eqref{eq:prediction-ball} with a continuous section. A sufficient choice
is $r<r_{\rm gate}$ and $Lr\le\sigma/2$, where $L$ bounds the Hessian of
$\phi$ on that gate slice. The map $p(u)$ is a ratio of finite polynomial
moment expressions; the prefix evidence is at least $\eta^n$ there.
Consequently its derivatives and such an $L$ are bounded directly by
finite differentiation using the calibration matrix and prior moments
through degree $q(n+m)$. This supplies checkable, generally nonoptimal
constants. It does not assert a uniform radius as $n,m$ grow.

\begin{theorem}[Attainable sharp memory--risk rate]\label{thm:memory-rate}
Let $\mathcal R_M^*$ be the infimum over all encoders with at most $M\ge1$
stored symbols and all decoders of their worst-case conditional regret in
the preceding fixed score problem, over attainable histories after $n$
reports. Encoders need not be continuous. Fresh private randomness is permitted. A shared public seed may be drawn
at the checkpoint independently of the entire past, the parameter and the
cartridges; it carries no historical information other than the $M$-symbol
message. Any retained seed that generated past commands counts as part of
that message. The query $J$ is independent of all these seeds and of the past. Then
\begin{equation}\label{eq:sharp-memory-rate}
 \frac{d}{d+2}\frac{\rho^2}{s}M^{-2/d}
 \ \le\ \mathcal R_M^*
 \ \le\ 4d\,M^{-2/d},\qquad d=q\min(n,m).
\end{equation}
In particular, with $B$ mutable bits, the optimal regret has order
$2^{-2B/d}$, with positive constants for the fixed apparatus, prior,
probe list and checkpoint. The smallest $B$ guaranteeing regret at most
$\varepsilon$ is
\begin{equation}\label{eq:bits-law}
 B_*(\varepsilon)=\frac d2\log_2(1/\varepsilon)+O(1)
 \quad (\varepsilon\downarrow0).
\end{equation}
No uniform lower constant over vanishing detector amplitude or all
full-support priors is asserted.
\end{theorem}
\begin{proof}
For the upper bound the whole prediction set lies in a $d$-dimensional
affine space and in $[0,1]^s$. Choose one point in it and orthonormal
coordinates in that affine space. Every coordinate then has absolute value
at most $\sqrt s$. Put $k=\lfloor M^{1/d}\rfloor$, and partition each
coordinate interval $[-\sqrt s,\sqrt s]$ into $k$ equal intervals. Encode
the nearest product-grid center. There are at most $k^d\le M$ centers,
and squared error in the ambient Euclidean norm is at most $ds/k^2$.
Clip the decoded probability coordinates to $[0,1]$; this cannot increase
the error to a true probability vector. Since $k\ge M^{1/d}/2$, the
regret is at most $4dM^{-2/d}$. This is an actual finite lookup encoder and
decoder on the exact moment or coefficient state, not an existence claim
about discontinuous real encodings.

For the lower bound use the uniform volume distribution on the radius-$\rho$
ball \eqref{eq:prediction-ball}, transported to histories by its section.
First fix any realization of public randomness. Conditional on a message
and $J$, averaging any decoder randomness can only reduce squared error.
Its at most $M$ conditional-mean output vectors are therefore code centers.
An encoder's private randomization cannot beat the closest center at a
fixed $p$. Orthogonally project the centers to the affine $d$-plane;
projection cannot increase distance to points of the ball. For a uniform
point $X$ of that ball and any such centers $z_1,\ldots,z_M$,
\[
 \Pr\{\min_i\|X-z_i\|\le r\}
 \le\min\{1,M(r/\rho)^d\}.
\]
This is the volume bound for a union of $M$ radius-$r$ balls, allowing
centers outside the attainable set. Integration of the complementary tail
from $0$ to $\rho^2M^{-2/d}$ gives
\[
 \mathbb E\min_i\|X-z_i\|^2
 \ge\int_0^{\rho^2M^{-2/d}}
       \left(1-\frac{M t^{d/2}}{\rho^d}\right)dt
 =\frac d{d+2}\rho^2M^{-2/d}.
\]
The bound is uniform in the fixed public seed, so averaging that seed
preserves it. Worst-case regret is at least this integrated regret divided
by $s$. This proves the lower bound for all declared randomized codes.
Putting $M=2^B$ and solving the two inequalities proves
\eqref{eq:bits-law}. These elementary covering arguments use the classical
squared-error quantization mechanism \cite{GrafLuschgy}; the content
specific to this experiment is the proven attainable decision dimension
and the common physically executed score task.
\end{proof}

\begin{theorem}[The same exponent in one fixed randomized experiment]\label{thm:average-memory-rate}
Fix the apparatus and the score query protocol above. Before each of the
first $n$ cartridges draw the finite command vector independently and
uniformly from its lookup cube, independently of the radius and all
cartridges. The commanded values and reports are known to the encoder.
At the checkpoint their randomization seed, commands and reports are
available to the decoder only through the $M$-symbol message. The
exploration rule itself is fixed and known. Let $\overline{\mathcal R}_M^*$
be the smallest unconditional expected regret, compared with the
full-history Bayes forecaster, over all such codes. Then for fixed $n,m$
there are $c_*>0$ and $C_*<\infty$ such that
\begin{equation}\label{eq:average-memory-law}
 c_*M^{-2/d}\le\overline{\mathcal R}_M^*\le C_*M^{-2/d},
 \qquad d=q\min(n,m).
\end{equation}
One may take $C_*=4d$. The identical independent uniform four-entry
exploration rule, followed by the shared query menu, gives this law in
both sign experiments, with $d=3\min(n,m)$ and $d=2\min(n,m)$ respectively.
Redundant command entries in the erased device are averaged as before.
\end{theorem}
\begin{proof}
Let $K$ be the number of scalar command entries in the complete prefix,
and write $g$ for that vector. Its fixed distribution has constant density
$(1-2\eta)^{-K}$ on its cube. On the all-failure event the evidence is
$Z(g)=\pi_0\prod_iF_{g_i}\ge\eta^n$, and its prediction vector is the
smooth map
\[
 p(g)=\left(\frac{\pi_0(H_j\prod_iF_{g_i})}{Z(g)}\right)_{j=1}^s.
\]
The product and Gram arguments give an interior $g_0$ at which this map
has rank $d$. Choose an orthonormal basis $E$ for its affine tangent
space and $K-d$ complementary command coordinates $w$. The map
\[
 g\longmapsto\bigl(E^\top(p(g)-p(g_0)),w(g)\bigr)
\]
has invertible derivative at $g_0$. By the inverse function theorem it has
a smooth inverse $\Phi(y,w)$ on a neighborhood. Choose a closed product
$B_\rho^d\times W$ inside that neighborhood, with $W$ a positive-volume
box in the complementary coordinates, such that its inverse image remains
strictly inside the command cube. Compactness gives
$|\det D\Phi|\ge j_0>0$ there. With the usual unit volume convention in
zero complementary dimensions, the same formula covers $K=d$.

The joint measure of $g$ and an all-failure prefix has density
$(1-2\eta)^{-K}Z(g)$. Changing coordinates and integrating over $W$ shows
that the subprobability distribution of $p$ on this prefix dominates
$\beta$ times uniform volume on $p(g_0)+EB_\rho^d$, where one can take
\[
 \beta=\frac{\eta^n j_0\,|W|\,|B_\rho^d|}{(1-2\eta)^K}>0.
\]
This is an unconditional subprobability minorization, not a posterior
bound obtained by ignoring the probability of failure. It also shows
$\beta\le1$. At most $M$ conditional-mean decoded vectors are available
after fixing any fresh public coding seed. No encoder, even one that uses
the complementary coordinates or private randomness, can beat distance
to the nearest of these centers. The volume calculation from
Theorem~\ref{thm:memory-rate} therefore gives expected regret at least
\[
 \beta\frac{d}{d+2}\frac{\rho^2}{s}M^{-2/d}.
\]
Regret is nonnegative on all other histories. Averaging the fresh coding
seed preserves this lower bound. The same universal grid code as before
has regret at most $4dM^{-2/d}$ on every history, proving the upper bound.

For the two apparatuses use exactly the same uniform four-entry commands.
The erased map has redundant coordinates, but its attainable coefficient
map still has rank three, so its prediction differential has rank
$2\min(n,m)$. The same argument with its own complementary coordinates
applies to that identical command law. The collision-sign device has rank
$3\min(n,m)$. The query and cartridge counts are unchanged.
\end{proof}

Here the randomness that generated past commands is part of the prefix,
not the fresh public coding seed. The continuous command distribution is
an explicit part of this fixed experiment. With a fixed finite command
alphabet and a finite report alphabet there are only finitely many
prefixes, and an exact finite code eventually suffices; no contrary
asymptotic law is asserted for that different input protocol. The result
above supplies a single fixed randomized experiment with a nondegenerate
attainable prediction distribution, rather than relying only on a
worst-command conditional family.

\begin{theorem}[Continuous real memory and a positive regret threshold]\label{thm:continuous-regret}
A deterministic state that is continuous in known command histories and
has $k<d$ real coordinates cannot obtain worst-case regret below
$\rho^2/s$ in the same score task, regardless of the decoder's measurability
or its private randomization. A continuous $d$-coordinate state obtains
zero regret. This assertion is separate from the finite-alphabet result;
it does not impose continuity on finite-bit encoders or allow a
history-dependent uncharged public seed.
\end{theorem}
\begin{proof}
Restrict the section in \eqref{eq:prediction-ball} to its boundary sphere.
The encoder composed with it is a continuous map $S^{d-1}\to\mathbb R^k$.
The Borsuk--Ulam theorem gives antipodal points with the same state
\cite{Matousek}. For $d=1$ the target is a point and this conclusion is
immediate. Their predictions are $p_*\pm\rho Eu$. For any common decoded
vector $a$, the average of its two squared distances is
$\|a-p_*\|^2+\rho^2\ge\rho^2$. The same identity holds after averaging
private decoder randomness. At least one regret is therefore at least
$\rho^2/s$. The exact $d$-coordinate prediction state gives the matching
zero-regret construction.
\end{proof}

\subsection{Identical-budget sign erasure changes the sharp law}
Use the same prior, the same $n,m$, the same four-entry command convention,
the same $4^m$ public query words, and the same observable score for the two
apparatuses in Section~\ref{sec:calibrated}. With the sign available,
\[
 d=3\min(n,m),\qquad
 B_*(\varepsilon)=\tfrac32\min(n,m)\log_2(1/\varepsilon)+O(1).
\]
With the sign erased before command evaluation and reporting,
\[
 d=2\min(n,m),\qquad
 B_*(\varepsilon)=\min(n,m)\log_2(1/\varepsilon)+O(1).
\]
Both statements hold for worst-history regret and, by
Theorem~\ref{thm:average-memory-rate}, unconditional regret under the same
fixed random exploration. They follow from the explicitly computed
attainable spans. Unlike a report-cost-saving example, removing the
structural feature changes the sharp exponent for the very same score
protocol. This is a law of attainable forecasting accuracy versus memory;
it is not a claim that a richer observation is worse when memory is free.
The next section supplies its positive information-value counterpart.

The lower bounds concern retained histories of positive probability. For
any fixed commanded all-failure prefix the probability of that event is at
least $\eta^n$. Comparing a code with the full-history Bayes forecaster,
Brier regret is nonnegative on every other history. Hence the corresponding
worst-command unconditional expected regret is at least $\eta^n$ times the
conditional lower bound. This makes explicit the price of the checkpoint
conditioning; it neither charges failed cartridges zero nor asserts a
horizon-uniform constant. The minimax upper bound holds on all histories.

\paragraph{Perturbations and computation.}
Invertibility of the finite calibration matrix is open: if its smallest
singular value is $\sigma>0$ and a same-degree perturbation has operator
norm below $\sigma$, its rank persists. Positivity and normalization must
also be retained. Under these hypotheses the attainable-ball and Gram
arguments apply anew and give the same exponents, with new constants.
For a nonpolynomial perturbation no exact rank conclusion follows from
small total variation; Section~\ref{sec:robust} instead controls common-policy
values. No rare-history posterior estimate or asymptotically zero error
floor is inferred from that value comparison. The finite-memory theorem
charges only history storage at the checkpoint, not the computation of an
optimal codebook, the query table size, or global control optimization.
